\chapter{人造卫星与行星运动}\label{ch:g12:satellites-kepler}

阳台上的碟形天线被牢牢固定，却整天凝视天空中的一个固定点，
该点高出地面 \qty{36000}{km}；空间站只需几分钟便掠过傍晚的天空；
由飞行“时钟”组成的星座，
告诉你的手机你身处何地。这些装置都没有
开动发动机：它们全都只是在\emph{下落}。本章
将计算每个装置应以多快的速度、在多高处下落——然后
用同一定律“称量”行星。

\section{绕地球的下落}

一切都建立在前面所测得的定律之上
（见\cref{ch:g10:universal-gravitation}）：质量为 $M$ 的地球
对质量为 $m$ 的卫星施加引力，二者中心相距 $r$，
引力方向指向地心，大小为
\[ F = G\,\frac{m M}{r^2}, \qquad G = \qty{6.67e-11}{N.m^2/kg^2}. \]
注意 $r$ 的含义：它是两物体\emph{中心之间}的距离。若卫星
离地高度为 $h$，则 $r = R + h$，其中 $R = \qty{6.37e6}{m}$ 为地球
半径——忘记 $R$ 是本章最典型的错误。

\begin{remark}[牛顿大炮，再谈]\label{rem:g12:satellites-kepler:cannon}
水平发射炮弹，炮弹便会落地；速度越大，落点越远；
速度恰当时，地面弯离的速度恰与
炮弹下落的速度相同，下落轨迹自行闭合，形成轨道。卫星
无需发动机，只需恰当的\emph{侧向}速度。
\end{remark}

\begin{omfigure}
\begin{tikzpicture}[scale=1.0, font=\small]
  \fill[omExo!15] (0,0) circle (1.5); \draw[omExo, thick] (0,0) circle (1.5);
  \fill[omExo!60] (-0.22,1.44) -- (0,1.85) -- (0.22,1.44) -- cycle;
  \coordinate (L) at (0,1.85);
  \draw[omProp, thick] (L) .. controls (0.7,1.85) .. (50:1.5);
  \draw[omProp, thick] (L) .. controls (1.4,1.85) and (2.0,0.9) .. (0:1.5);
  \draw[-Stealth, omDef, thick] (0,1.85) arc (90:-260:1.85);
  \draw[-Stealth, omThm, thick] (L) .. controls (1.7,2.0) and (3.2,2.2)
    .. (4.3,2.9);
  \node[omDef, anchor=east] at (-1.85,1.1) {轨道};
  \node[omThm, anchor=west] at (3.35,2.35) {逃逸};
  \node[omProp, anchor=west] at (1.35,0.55) {下落};
\end{tikzpicture}

{\small 牛顿大炮：发射速度越大，炮弹落点越远（橙色）；
速度约为 \qty{7.9}{km/s} 时，其下落与地球表面的弯曲恰好吻合，炮弹绕地球运行（蓝色）；
速度更大时，炮弹便会逃逸（红色）。}
\end{omfigure}

\section{圆轨道：一个速率，一个周期}

\begin{theorem}[轨道速率]\label{thm:g12:satellites-kepler:speed}
卫星在半径为 $r$ 的匀速圆轨道上绕
质量为 $M$ 的天体运行，其速率为
\[ v = \sqrt{\frac{G M}{r}}, \]
它与卫星的质量无关。
\end{theorem}

\begin{proof}
重力是唯一的力，故由牛顿第二定律
（见\cref{thm:g12:newtons-laws:second}）得 $m\vect a = \vect F$；
将其投影到弗雷内标架上（见\cref{def:g12:kinematics-2d:frenet}）。
引力指向圆心，其切向分量为零，故速率保持不变；
其法向分量即为全部 $GmM/r^2$，
匀速圆周运动的法向加速度为 $a_n = v^2/r$
（见\cref{thm:g12:kinematics-2d:centripetal}）。由此 $m v^2/r =
GmM/r^2$，并——约去 $m$——得 $v^2 = GM/r$。
\end{proof}

\begin{omfigure}
\begin{tikzpicture}[scale=1.0, font=\small]
  \fill[omExo!15] (0,0) circle (0.85); \draw[omExo, thick] (0,0) circle (0.85);
  \node at (0,0) {地球}; \draw[dashed, omExo] (0,0) circle (2.3);
  \draw[thin] (0,0) -- (40:2.3) node[midway, above left] {$r$};
  \fill (40:2.3) circle (2.2pt) node[above right] {$S$};
  \draw[-Stealth, omThm, very thick] (40:2.3) -- (40:1.35)
    node[pos=0.85, below right] {$\vect F$};
  \draw[-Stealth, omDef, thick] (40:2.3) -- ++(130:1.15)
    node[above] {$\vect v$};
\end{tikzpicture}

{\small 轨道运动的全部机理：速度沿切线方向，力
指向圆心——引力从不减慢卫星，它只是使其轨迹弯曲，
并且持续不断。}
\end{omfigure}

\begin{proposition}[轨道周期]\label{prop:g12:satellites-kepler:period}
卫星绕行一周所需的周期为
\[ T = \frac{2\pi r}{v} = 2\pi \sqrt{\frac{r^3}{G M}}. \]
\end{proposition}

\begin{proof}
绕行一周的路程为圆周长 $2\pi r$，速率恒定；代入
$v = \sqrt{GM/r}$，得 $T = 2\pi r\sqrt{r/(GM)} = 2\pi\sqrt{r^3/(GM)}$。
\end{proof}

$v$ 和 $T$ 都不含 $m$：从空间站上意外松脱的螺钉，
会与空间站一同沿轨道运行。二者都与 $r$ 有关：
\emph{轨道越高，卫星运行越慢}，无论是速率
（$v \propto 1/\sqrt r$），还是周期（$T \propto r^{3/2}$）。

\begin{example}[空间站]\label{ex:g12:satellites-kepler:iss}
国际空间站的飞行高度为
$h \approx \qty{400}{km}$，故 $r = \num{6.37e6} + \num{4.0e5} =
\qty{6.77e6}{m}$；取 $GM = \num{6.67e-11} \times \num{5.97e24} =
\qty{3.98e14}{m^3/s^2}$, $v = \sqrt{\num{3.98e14}/\num{6.77e6}}
\approx \qty{7.7e3}{m/s}$，而 $T = 2\pi r/v \approx \qty{5.5e3}{s}
\approx \qty{92}{min}$：每隔一个半小时绕地球一周，
宇航员每天约能看到十六次日出。
\end{example}

\begin{example}[月球像苹果一样下落]\label{ex:g12:satellites-kepler:moon}
月球位于 $r = \qty{3.84e8}{m}$ 处，绕地球运行的周期为 $T = \qty{27.3}{d} =
\qty{2.36e6}{s}$，故 $v = 2\pi r/T \approx \qty{1.02e3}{m/s}$，
$a = v^2/r \approx \qty{2.7e-3}{m/s^2}$——恰为 $g/60^2$，而
月球位于距地心 $60$ 个地球半径处：使苹果下落的那股拉力
同样的拉力使月球轨道弯曲，且按平方反比减弱。每一秒
月球朝我们“下落”约 \qty{1.4}{mm}，其
侧向前进的一千米带着它绕行——这正是牛顿在 1666 年所做的验算。
\end{example}

\section{开普勒定律}

牛顿并非凭空猜出平方反比关系，而是从三条规律中提取出它，
这三条规律是开普勒于 1609 至 1619 年间归纳出的，
依据第谷·布拉赫二十年间肉眼观测的行星位置——
这是望远镜出现以前最好的数据。

\begin{definition}[椭圆]\label{def:g12:satellites-kepler:ellipse}
\emph{椭圆}\index{椭圆}是这样一些点的集合：这些点到两个定点
（即\emph{焦点}\index{焦点（椭圆）}）的距离之和为常数；
其最长直径的一半即为
\emph{半长轴}\index{半长轴} $a$，圆则是
两个焦点重合的情形。行星轨道上离太阳最近和最远的点
分别称为\emph{近日点}\index{近日点}
和\emph{远日点}\index{远日点}。
\end{definition}

\begin{theorem}[开普勒定律]\label{thm:g12:satellites-kepler:kepler}
\index{开普勒定律}%
对绕太阳运行的行星（质量为 $M$）：
\begin{enumerate}
  \item 每条轨道都是椭圆，太阳位于其一个焦点上；
  \item 太阳与行星的连线在相等时间内扫过相等的面积；
  \item 周期平方与半长轴立方之比
        对所有行星都相同：$T^2/a^3 = 4\pi^2/(G M)$。
\end{enumerate}
同样的定律支配绕任意中心天体运行的任一族卫星，
其中 $M$ 为中心天体的质量。\admitted
\end{theorem}

\begin{remark}[本阶段能证明的内容]\label{rem:g12:satellites-kepler:proof}
对于\emph{圆周}轨道，第三定律是我们能证明的：将
\cref{prop:g12:satellites-kepler:period}两边平方可得
$T^2/r^3 = 4\pi^2/(GM)$——对同一中心天体的每颗卫星，这是同一个常数。
椭圆轨道、等面积定律以及一般的 $a$ 情形都可由
平方反比律推出，这需要更多微积分，将在本科一年级
分册中给出。第二定律可这样理解：行星
在近日点运动最快，那里的扫掠线段较短，须迅速掠过。
\end{remark}

\begin{omfigure}
\begin{tikzpicture}[scale=0.85, font=\small,
    declare function={re(\t)=2.025/(1+0.5*cos(\t));}]
  \draw[thick] plot[domain=0:360, samples=120, smooth]
    ({1.35+re(\x)*cos(\x)}, {re(\x)*sin(\x)});
  \fill[omDef!30] (1.35,0) -- plot[domain=-40:40, samples=25]
    ({1.35+re(\x)*cos(\x)}, {re(\x)*sin(\x)}) -- cycle;
  \fill[omDef!30] (1.35,0) -- plot[domain=175:185, samples=10]
    ({1.35+re(\x)*cos(\x)}, {re(\x)*sin(\x)}) -- cycle;
  \fill[omThm] (1.35,0) circle (2.6pt) node[below=2pt] {太阳};
  \draw[-Stealth, omDef, thick]
    ({1.35+re(40)*cos(40)}, {re(40)*sin(40)}) -- ++(-0.55,0.42);
  \draw[-Stealth, omDef, thick]
    ({1.35+re(175)*cos(175)}, {re(175)*sin(175)}) -- ++(-0.12,-0.5);
  \node[right] at (2.75,0) {$P$}; \node[left] at (-2.72,0) {$A$};
\end{tikzpicture}

{\small 开普勒第二定律：在近日点 $P$ 附近的一个月内，行星
扫过一个短而宽的扇形；在远日点 $A$ 附近，扫过一个长而窄的
扇形——面积相等：离太阳近时行进快，离太阳远时行进慢。}
\end{omfigure}

\begin{omfigure}
\begin{tikzpicture}
\begin{axis}[omaxis, width=8.8cm, height=5.8cm,
    xmode=log, ymode=log,
    xlabel={半长轴 $a$（天文单位）},
    ylabel={周期 $T$（年）},
    xmin=0.25, xmax=50, ymin=0.15, ymax=400]
  \addplot[omThm, thick] coordinates {(0.3,0.1643) (40,253)};
  \addplot[only marks, mark=*, mark size=1.7pt, omDef] coordinates
    {(0.387,0.241) (0.723,0.615) (1,1) (1.524,1.881) (5.203,11.86)
     (9.537,29.46) (19.19,84.02) (30.07,164.8)};
  \node[below right, font=\scriptsize] at (axis cs:0.387,0.241) {水星};
  \node[below right, font=\scriptsize] at (axis cs:1,1) {地球};
  \node[below right, font=\scriptsize] at (axis cs:5.203,11.86) {木星};
  \node[below right, font=\scriptsize] at (axis cs:30.07,164.8) {海王星};
\end{axis}
\end{tikzpicture}

{\small 在双对数坐标轴上的八颗行星：斜率为
$3/2$，即 $T^2 \propto a^3$——开普勒第三定律，它在
这些单位（天文单位、年）下简化为 $T^2 = a^3$。}
\end{omfigure}

\section{轨道的应用}

\begin{definition}[地球静止轨道]\label{def:g12:satellites-kepler:geo}
若卫星悬停于地面某一固定点的上空，则称其为
\emph{地球静止}\index{地球静止轨道}。其轨道须为圆形、位于赤道面内，
且周期为一个\emph{恒星日}\index{恒星日}，
$T = \qty{86164}{s}$（\qty{23}{h}
\qty{56}{min} \qty{4}{s}）：地球自转一周所需的时间，
是\emph{相对于恒星}而言的。（\qty{24}{h}的一日是相对于
太阳而言的，地球每天也朝向太阳前进约一度。）
\end{definition}

\begin{example}[地球静止轨道的高度]\label{ex:g12:satellites-kepler:geoalt}
将开普勒第三定律反过来用，即可确定半径：
$r = \left(GMT^2/4\pi^2\right)^{1/3} =
\left(\num{3.98e14} \times (\num{86164})^2/4\pi^2\right)^{1/3}
\approx \qty{4.22e7}{m}$——高度 $h = r - R \approx
\qty{3.58e7}{m} \approx \qty{35800}{km}$，运行速率为 $v = 2\pi r/T
\approx \qty{3.1}{km/s}$。每一颗看似“静止不动”的中继卫星都位于
赤道上空的这条圆周上——别无他处。
\end{example}

\begin{example}[卫星导航星座]\label{ex:g12:satellites-kepler:gps}
卫星导航系统使约三十颗卫星运行于
$r \approx \qty{2.66e7}{m}$（高度 $\approx \qty{20200}{km}$），
由\cref{prop:g12:satellites-kepler:period}可得 $T \approx
\qty{4.32e4}{s}$：约半个恒星日，因此每颗卫星每天
重复一次其地面轨迹。你的接收机比较其时钟信号的到达时刻，
并与精确到米的已知轨道作比较。
\end{example}

\begin{method}[由卫星称量天体的质量]\label{met:g12:satellites-kepler:weigh}
\index{称量行星}%
要测量某颗行星或恒星的质量：找到它的任一颗卫星
（月球、探测器、行星均可），测出其轨道半径 $r$ 和周期
$T$，再由第三定律反解：
\[ M = \frac{4\pi^2 r^3}{G\,T^2}. \]
这只能称量\emph{中心}天体的质量——卫星自身的质量
在\cref{thm:g12:satellites-kepler:speed}中已消去。数据卡上的每个
质量都是这样得到的：地球由月球得出，太阳由
地球自身的轨道得出。
\end{method}

\begin{remark}[轨道越高越慢——但代价更大]\label{rem:g12:satellites-kepler:energy}
由 $v = \sqrt{GM/r}$ 可知，月球的速率仅为 \qty{1}{km/s}，而
空间站则以 \qty{7.7}{km/s} 疾驰。然而，较高的轨道
却\emph{代价更大}：逆着引力向上攀升，会储存势能
（账目在\cref{ch:g12:work-and-energy}中结清），而
$E_p$ 的增加量超过 $E_k$ 的减少量。因此，火箭需
\emph{向前}喷火才能进入更高、更慢的轨道；而
受稀薄空气阻滞的卫星则会\emph{向下}螺旋，并且速率增大。
\end{remark}

\begin{remark}[数据卡]\label{rem:g12:satellites-kepler:data}
除非另有说明，取 $G = \qty{6.67e-11}{N.m^2/kg^2}$；地球：
$M = \qty{5.97e24}{kg}$，$R = \qty{6.37e6}{m}$，恒星日
\qty{86164}{s}；月球轨道：$r = \qty{3.84e8}{m}$，$T =
\qty{27.3}{d}$；太阳：$M = \qty{1.99e30}{kg}$；日地距离
\qty{1.496e11}{m}；一年 $= \qty{3.156e7}{s}$；火星：$M =
\qty{6.42e23}{kg}$，$R = \qty{3.39e6}{m}$。
\end{remark}

\section{练习}

\begin{exercise}[$\star$]\label{exo:g12:satellites-kepler:1}
质量为 \qty{1000}{kg} 的卫星停在发射台上，之后在
高度 \qty{400}{km} 处飞行。分别计算两处地球对它的拉力：
飞行时其重量仍保留多少百分比？乘员为何会漂浮？
\end{exercise}

\begin{exercise}[$\star$]\label{exo:g12:satellites-kepler:2}
对一颗（理论上）掠过地球表面的卫星
（$r = R$），求其环绕速率和周期，并与
牛顿的炮弹实验（见\cref{rem:g12:satellites-kepler:cannon}）作比较。
\end{exercise}

\begin{exercise}[$\star$]\label{exo:g12:satellites-kepler:3}
将一颗卫星从半径为 $r$ 的圆轨道
移至半径为 $2r$ 的圆轨道。其速率和周期分别变为原来的多少倍？
答案是否与其质量有关？
\end{exercise}

\begin{exercise}[$\star$]\label{exo:g12:satellites-kepler:4}
重新计算空间站的数据：由 $r = \qty{6.77e6}{m}$ 出发，求
$v$、$T$，以及在 \qty{24}{h} 内完成的轨道圈数。
\end{exercise}

\begin{exercise}[$\star$]\label{exo:g12:satellites-kepler:5}
叙述开普勒三定律。彗星沿其极为扁长的椭圆运行时，
在何处速率最大？哪条定律说明了这一点？卫星的
周期是否与其自身质量有关？
\end{exercise}

\begin{exercise}[$\star\star$]\label{exo:g12:satellites-kepler:6}
称量地球：由月球的轨道半径和周期（数据
卡），求地球的质量，并与数据卡上的值作比较。
\end{exercise}

\begin{exercise}[$\star\star$]\label{exo:g12:satellites-kepler:7}
某公司提议发射一颗卫星，使其永久悬停在巴黎上空
（北纬 $49^\circ$）。根据重力的方向，说明
为何每条圆轨道的圆心都必须是地心；由此
断定该提议不可行。卫星\emph{能}悬停在何处？
\end{exercise}

\begin{exercise}[$\star\star$]\label{exo:g12:satellites-kepler:8}
火卫一绕火星运行，轨道半径 $r = \qty{9.38e6}{m}$，周期 $T = \qty{7}{h}
\qty{39}{min}$。推算火星的质量，并与数据卡作比较。
\end{exercise}

\begin{exercise}[$\star\star$]\label{exo:g12:satellites-kepler:9}
火星绕太阳运行，半长轴 $a = \qty{2.279e11}{m}$，周期为 \qty{687}{d}。
计算火星和地球（数据卡）的 $T^2/a^3$，检验
开普勒第三定律，并推算太阳的质量。
\end{exercise}

\begin{exercise}[$\star\star$]\label{exo:g12:satellites-kepler:10}
牛顿的月球检验：由数据卡求月球的速率及
其向心加速度 $a = v^2/r$，证明其等于 $g/60^2$，
并说明为何这使牛顿确信同一条定律支配苹果与月球。
\end{exercise}

\begin{exercise}[$\star\star$]\label{exo:g12:satellites-kepler:11}
以太阳系单位（天文单位、年）表示，开普勒第三定律为
$T^2 = a^3$，即本课程对数–对数图上的那条直线。在图上标出
小行星谷神星（$a = \qty{2.77}{au}$）与哈雷彗星
（$a = \qty{17.8}{au}$），并计算两者的周期。哈雷彗星上一次经过
近日点是在 1986 年——它何时回归？
\end{exercise}

\begin{exercise}[$\star\star\star$]\label{exo:g12:satellites-kepler:12}
稀薄空气擦过空间站，但其速率却随之\emph{增大}，
高度则在下降。分别计算 $v$ 在高度 \qty{400}{km} 和
\qty{350}{km} 处的值，再解释这一悖论：摩擦力使物体减速，
——那额外的速率从何而来？
\end{exercise}

\begin{exercise}[$\star\star\star$]\label{exo:g12:satellites-kepler:13}
一架望远镜发现了一颗系外行星绕其恒星运行，
轨道半径 $r = \qty{7.48e9}{m}$，周期为 \qty{3.50}{d}。求该恒星的质量，
分别以千克和太阳质量为单位表示。哪一个质量——恒星的还是行星的——
\emph{无法}用这种方法求得，为什么？
\end{exercise}

\begin{exercise}[$\star\star\star$]\label{exo:g12:satellites-kepler:14}
火星定居者希望在基地上空设置一个“火星静止”中继站，
火星相对恒星自转一周需 \qty{24}{h} \qty{37}{min}：
求该中继站的轨道半径，以及其离火星
表面的高度（数据卡）。
\end{exercise}

\begin{exercise}[$\star\star\star$]\label{exo:g12:satellites-kepler:15}
哈雷彗星经过近日点时，距太阳 $r_p = \qty{8.77e10}{m}$，
速率为 \qty{54.5}{km/s}；远日点距太阳 $r_a = \qty{5.25e12}{m}$。
在两个极端位置，速度均与太阳–彗星连线段垂直，
故由开普勒第二定律必有 $r_p v_p = r_a v_a$
（扫过的三角形面积相等，均为 $\tfrac12 r v\,\Delta t$）。由此
求远日点速率，并说明该比值的意义。
\end{exercise}

\section{周末问题：赤道上空的一个位置}

\begin{problem}[{周末问题——把卫星放到赤道上空：
为什么天空中唯一的“停车位”是一个圆，
位于离地 \qty{35800}{km} 处，它与下方高速运行的空间站相比如何，
以及同一定律如何顺带称出木星的质量}]
\label{pb:g12:satellites-kepler:1}
某电信公司需要一颗中继卫星，使其地面天线一经固定便永不必再追踪：
这就是地球静止卫星。你是任务分析员，
请使用数据卡。在第四部分中，木星的卫星木卫一
绕木星运行的轨道半径为 $r = \qty{4.22e8}{m}$，周期为 $T = \qty{1.77}{d}$，木星
相对恒星自转一周需要 \qty{9}{h} \qty{56}{min}。

\medskip
\noindent\textbf{第一部分——静止悬停的规则。}
\begin{enumerate}
  \item 卫星必须始终位于地面某一固定点的上空。它的周期必须为多少——
        以及为什么恒星日为
        \qty{86164}{s}，而非 \qty{24}{h}？
  \item 圆周运动中，合力指向圆心；此处唯一的力是引力，
        由此推出任一圆轨道的
        圆心都是地球中心。
  \item 若一个悬停在巴黎纬度上空的圆轨道，其圆心将位于地轴上，
        且位于地心以北。由此得出：
        地球静止轨道须位于哪一平面内？
  \item 在弗勒内标架中，对圆轨道写出牛顿第二定律，
        设轨道半径为 $r$，并导出 $v = \sqrt{GM/r}$。
  \item 由此推出 $T = 2\pi\sqrt{r^3/(GM)}$，并将其改写为开普勒
        第三定律 $T^2/r^3 = 4\pi^2/(GM)$。
\end{enumerate}

\medskip
\noindent\textbf{第二部分——唯一可行的圆轨道。}
\begin{enumerate}[resume]
  \item 由第三定律解出 $r$，并求地球静止轨道的
        半径。
  \item 由此求出离地面的高度。
  \item 由 $v = 2\pi r/T$ 计算轨道速率。
  \item 将结果与 $v = \sqrt{GM/r}$ 相比较以作检验。
  \item 该公司的卫星质量为 \qty{4.0e3}{kg}；竞争对手的
        卫星质量是其两倍。比较两者的地球静止轨道半径，
        并说明额外质量在发射时会带来什么影响。
\end{enumerate}

\medskip
\noindent\textbf{第三部分——下方的空间站。}
\begin{enumerate}[resume]
  \item 空间站的飞行轨道半径为 $r = \qty{6.77e6}{m}$：求其
        速率。
  \item 求其周期，分别以秒和分钟为单位。
  \item 空间站在 \qty{24}{h} 内绕地球运行多少圈——
        宇航员每天大约能看到多少次日出？
  \item 求半径之比 $r_{\text{地球静止}}/r_{\text{空间站}}$，
        并验证周期之比等于该数的
        $3/2$ 次幂，正如开普勒第三定律所要求。
  \item 哪一颗卫星运动得更快，哪一颗送入轨道时每千克
        所耗能量更多？请用一句话调和这两个结论。
\end{enumerate}

\medskip
\noindent\textbf{第四部分——顺带称量木星。}
\begin{enumerate}[resume]
  \item 对木卫一应用开普勒第三定律，将木星的
        质量 $M_J$ 用 $G$、$r$ 和 $T$ 表示。
  \item 计算 $M_J$。
  \item 这相当于多少个地球质量？大约是太阳质量的
        几分之一？
  \item 一个“木星静止”中继站应悬停在木星
        云层的某一点上方：求其轨道半径，并与
        木卫一的轨道半径相比较。
  \item 向董事会汇报两句话：公司应将中继卫星停放的高度，
        以及同一定律“免费”给出的
        木星质量。
\end{enumerate}
\end{problem}
